\scanpage{69}
\mspage{62}
% archive sheet number: 68
(pas seulement sur $M(I)_f$) $p_*$ est compatible avec localiseurs $W(I)$, $W$.

Si $I$ n'a pas de plus grand élém\unsure{t|.}, alors l'ens. \inserted{$M$} de ses éléments maximaux est de card $\geq2$, on peut donc l'écrire $M=M_1^{\unsure{\circ|o|'}}\amalg M_2^\circ$, et on aura
\[
I=\unsure{I_1|\mathbb I_1}\cup I_2,\quad\text{où }I_1=I_{/M_1^\circ},\quad I_2=I_{/M_2^\circ}
\]
sont deux ouverts de $I$, \unsure{différents|distincts} de $I$. \inserted{donc de card $n_1,n_2<n$.} Considérons aussi $I_0=I_1\cap I_2$, c'est un ouvert de card. $<n$. Le diagramme
\diagram{
\begin{tikzcd}[ampersand replacement=\&,column sep=large,row sep=small]
\& I_1\arrow[dr,dotted] \& \\
I_0\arrow[ur]\arrow[dr] \& \& I=I_1\cup I_2 \\
\& I_2\arrow[ur,dotted] \&
\end{tikzcd}
}
correspond à une flèche
\diagram{
$u:I\longrightarrow\unsure{\Psi|\Upsilon}={}$
\begin{tikzcd}[ampersand replacement=\&,row sep=small]
\&1\\
0\arrow[ur]\arrow[dr]\&\\
\&2
\end{tikzcd},
}
avec
\[
I_1=I_{/1},\quad I_2=I_{/2},\quad I_0=I_{/0}.
\]
L'hyp. de récurrence et a) implique que $u$ est admissible. Comme un composé de foncteurs admissibles est \inserted{admissible} (tautologiquement, par $(vu)_*\simeq v_*u_*$), on est réduit finalement au
